\documentclass{article}
\usepackage{amsmath,hyperref}
\begin{document}
\section*{Select the supplier at the pose-driven dispatch cut}
\paragraph{What the source makes different.}
In the pinned SubT sources, an incoming message is queued by
\texttt{Broker::OnMessage}. A pose event installs the advertised
positions, rebuilds the visibility table for newly static
breadcrumbs, and then dispatches the pending messages. Enqueueing
is not a reservation of the route seen at enqueue time. The fixture
executes those source bodies, instantiating pose/time conversion and
broker ownership as serial shims. Neighbor reporting is stubbed.
A proposed local selector uses that pre-dispatch point after the
visibility update; it is not an existing stock hook. It samples
the pinned deterministic range/visibility and exact live byte
counters at the same serial cut as effect. It cannot promise the
same result under concurrent mutation or stochastic radio.

\paragraph{Matched duties and controls.}
The old suffix is a 164-byte stand-in for the compact continuation
of PR 860. Both P and Q are assumed armed with the same exact copy;
the early byte has already arrived from P in every continuation.
This experiment tests the byte transport and dispatch order, not
copy authenticity, sender preparation or warehouse installation.
A separate four-byte job is still owed from P. Its route is
independent of the old report but the jobs share P's transmit budget.
Vary P's earlier debit $d=40,41,42$, Q's $q=150,151,152$ and the
new native pose event: unchanged world, static near, static far.

In a P-only cut, the old continuation is possible exactly when
$d+115+164\le320$. Paying the separate four bytes before it makes
it impossible throughout the tested band. In a Q-only cut, the old
continuation is possible exactly when $q+5+164\le320$, and the
auxiliary work can proceed if P is not charged an attempted old
transfer. A failed old attempt still uses P's transmit budget.
If both cuts remain possible after the same qualified prefix,
no fixed complete queue order both preserves the old suffix
whenever physically affordable and preserves auxiliary work wherever
it can coexist with the old suffix:
putting the P auxiliary job first sacrifices the old P-only case;
putting the P suffix first loses the auxiliary job even in a
Q-only case where Q can pay for and deliver the old suffix.
Reversing the two redundant suffixes without canceling the old
P work cannot remove that debit. This obstruction is for dispatching
the scheduled suffixes blindly with the stated dedicated send budgets,
not a theorem excluding other codecs or qualified feedback.

\paragraph{Selection is not allocation.}
The local intervention can select one live sender and meter just
that attempt. Preserving the old transfer first does not determine
the best supplier when both fit: Q leaves P's capacity available
for the independent work. The case sweep checks delivery and every
sender and receiver counter. A direct same-cut policy with the
same access chooses Q whenever feasible, then P, and attains the
same bound. A serial prototype at the source ordering point is
evidence for earned coordination access, not new physical authority,
a stock production capability, or a full independent mission.
\end{document}
